\section{Conclusion}
\label{sec:conclusion}

This project demonstrated the application of Integer Linear Programming to three logic puzzles of contrasting difficulty.

\textbf{Mosaic} (difficulty 1) is an ideal ILP problem: all constraints are purely local, the formulation is compact, and CPLEX solves every tested instance in well under one second regardless of grid size or clue density. The constraint propagation heuristic, while fast on small grids, is unreliable on larger sparse instances due to the absence of backtracking.

\textbf{Range} (difficulty 5) presents a significantly richer modelling challenge. The visibility constraints required a careful linearisation through auxiliary binary variables, and white-cell connectivity required a global enforcement mechanism. Two strategies were implemented and compared:
\begin{itemize}
  \item The \textbf{single-commodity flow} formulation encodes connectivity statically but at the cost of Big-M coefficients that weaken the LP relaxation.
  \item The \textbf{lazy-constraint callback} formulation avoids any extra variables upfront and enforces connectivity dynamically via separator inequalities added as CPLEX discovers integer-infeasible candidates. This yields a tighter LP relaxation and consistent 2--6$\times$ speed improvements across all instance classes.
\end{itemize}
Both CPLEX approaches solve every generated instance optimally within seconds. The Range heuristic, based on local propagation of visibility rays, is largely ineffective on random instances because Range's constraints are inherently global and non-local.

\textbf{Galaxies} (difficulty 9) is the most complex puzzle tackled in this project. Rotational symmetry is encoded directly as linear equality constraints, while region connectivity --- which cannot be compactly encoded --- is handled exclusively via lazy-constraint callbacks. The callback submits the standard separator cut $\sum_{S}(1{-}x) + \sum_{N(S)} x \geq 1$ on every disconnected component; an optional symmetric-twin cut on $\sigma_k(S)$ is also implemented but, on our 25-instance dataset, did not measurably accelerate convergence (the underlying constraints are tight enough that CPLEX rediscovers the mirror configuration cheaply on its own). CPLEX solves all 25 instances (sizes $4 \times 4$ to $15 \times 15$), with mean solve times below $10$\,ms on grids up to $7 \times 7$ and a worst-case of $\sim 2$\,s on adversarial $10 \times 10$ instances. The propagation-based heuristic succeeds on $\sim 43\,\%$ of instances; on this dataset, providing it as a CPLEX warm start did not measurably change solve times either way.

\paragraph{Cross-puzzle synthesis.} The three puzzles span an interesting axis of constraint locality. Mosaic's constraints are purely local (each binary appears in at most nine clue equations); CPLEX solves the LP relaxation at the root and the gap to integer-feasibility is essentially zero. Range adds a global \emph{connectivity} requirement, which our two formulations make explicit either statically (flow with Big-M) or dynamically (lazy cuts); the latter wins by orders of magnitude on small instances and only converges to comparable performance on the larger ones. Galaxies adds a structural \emph{symmetry} axis on top: cells gain a long-range constraint to their mirror images, but the partitioning is also strongly anchor-driven, which means the LP relaxation remains tight and only a small fraction of instances ever need the connectivity callback to fire. Across the three, the lesson is that ILP is at its best when constraints are local \emph{or} when the global constraint can be lazily separated by a polynomial-time oracle (BFS in our case).
